\chapter{Arquitectura de DBMS, Independencia de Datos y Transacciones ACID}
\label{cap:arquitectura_dbms_acid}

\section{Marco Conceptual y Te\'orico (Curr\'iculo Universitario)}

\subsection{Arquitectura ANSI/SPARC de Tres Niveles}
El est\'andar ANSI/SPARC define tres niveles de abstracci\'on esenciales para desacoplar las aplicaciones de la persistencia f\'isica:
\begin{enumerate}
    \item \textbf{Nivel Externo (Vistas de Usuario):} Describe la parte de la base de datos relevante para un grupo de usuarios particular mediante vistas y restricciones de acceso.
    \item \textbf{Nivel Conceptual (Esquema L\'ogico):} Modela la estructura l\'ogica global (tablas, llaves primarias, llaves for\'aneas y reglas de integridad de negocio).
    \item \textbf{Nivel Interno (Esquema F\'isico):} Define c\'omo se distribuyen los bloques en disco, estructuras de almacenamiento (Heap files, B-Trees) y optimizaci\'on de I/O.
\end{enumerate}

\subsection{Independencia de Datos}
\begin{itemize}
    \item \textbf{Independencia L\'ogica:} Capacidad de modificar el esquema conceptual (a\~nadir nuevas columnas o tablas) sin alterar las aplicaciones existentes ni las vistas externas.
    \item \textbf{Independencia F\'isica:} Capacidad de reorganizar el almacenamiento interno (a\~nadir \'indices, cambiar m\'etodos de acceso o particionar tablas) sin alterar el esquema conceptual.
\end{itemize}

\subsection{El Modelo Transaccional ACID}
\begin{itemize}
    \item \textbf{Atomicidad (A):} La transacci\'on se ejecuta como una unidad indivisible: o se completan todas las operaciones (\texttt{COMMIT}) o se revierten por completo (\texttt{ROLLBACK}).
    \item \textbf{Consistencia (C):} La transacci\'on transforma la base de datos de un estado v\'alido a otro, cumpliendo todas las restricciones de integridad.
    \item \textbf{Aislamiento (I):} Las transacciones concurrentes operan sin interferencias no controladas (Niveles: \textit{Read Uncommitted, Read Committed, Repeatable Read, Serializable}).
    \item \textbf{Durabilidad (D):} Una vez confirmado el commit, los cambios persisten incluso ante ca\'idas del servidor gracias al registro de transacciones (\textit{Write-Ahead Logging / WAL}).
\end{itemize}

---

\section{Casos de Estudio: Control Concurrente y Bloques Transaccionales}

\begin{lstlisting}[style=sqlstyle, caption={Manejo seguro de transacciones con bloqueo expl\'icito}]
-- Bloque transaccional seguro con actualizaci\'on at\'omica de inventario
BEGIN;

-- Bloqueo pesimista a nivel de fila para evitar condiciones de carrera (Race Conditions)
SELECT
    product_id,
    product_name,
    units_in_stock
FROM products
WHERE product_id = 11
FOR UPDATE;

-- Actualizaci\'on condicional segura
UPDATE products
SET units_in_stock = units_in_stock - 5
WHERE product_id = 11
  AND units_in_stock >= 5;

COMMIT;
\end{lstlisting}

---

\section{Taller de Consolidaci\'on del Cap\'itulo}

\begin{businessproblem}
Una plataforma de comercio electr\'onico necesita procesar la cancelaci\'on de un pedido de compra, reintegrando autom\'aticamente el stock f\'isico a las tablas de inventario y marcando la orden como anulada en una sola transacci\'on indivisible para evitar desajustes de auditor\'ia.
\end{businessproblem}

\subsection{Soluci\'on T\'ecnica en PostgreSQL}

\begin{lstlisting}[style=sqlstyle, caption={Transacci\'on ACID de anulaci\'on y restituci\'on de inventario}]
-- Transacci\'on segura con restituci\'on autom\'atica y RETURNING
BEGIN;

-- 1. Actualizar estado de la orden
UPDATE orders
SET
    shipped_date = NULL,
    ship_name = 'ANULADO - ' || ship_name
WHERE order_id = 10248
RETURNING order_id, customer_id;

-- 2. Restituir el stock de cada producto involucrado
UPDATE products AS p
SET units_in_stock = p.units_in_stock + od.quantity
FROM order_details AS od
WHERE od.order_id = 10248
  AND p.product_id = od.product_id;

COMMIT;
\end{lstlisting}

\begin{engtip}
\begin{itemize}
    \item \textbf{Integridad de Transacci\'on:} Al encerrar ambas sentencias entre \texttt{BEGIN} y \texttt{COMMIT}, si la base de datos se reinicia o falla durante el segundo \texttt{UPDATE}, PostgreSQL ejecuta un \texttt{ROLLBACK} autom\'atico protegiendo el balance contable.
    \item \textbf{Uso de \texttt{RETURNING}:} Elimina la necesidad de ejecutar un \texttt{SELECT} previo para confirmar el identificador procesado.
\end{itemize}
\end{engtip}
